\chapter{Самосогласованное внутреннее время и предел временной ловушки}

Ранняя постановка проекта содержала ещё один механизм устойчивости,
предшествовавший четырёхтактному резонатору. Предполагалось, что быстрый
вихрь создаёт собственную ``яму времени'': внутри частицы собственное время
идёт обычно, но относительно внешнего наблюдателя процесс замедляется.
В историческом тексте этому механизму приписывался безразмерный радиус
$R=1{,}33$, якобы полученный из системы Шрёдингера--Ньютона.

Настоящая глава отделяет три разных утверждения:
\begin{enumerate}
 \item кинематическое замедление внутреннего процесса;
 \item динамическую устойчивость ненулевого дефекта;
 \item возможность обмена с внешним световым каналом.
\end{enumerate}
Их совпадение нельзя предполагать заранее.

\section{Что действительно сохранилось в архиве}

В файле
\path{архив-2025-2026/2025-12-истоки/habr/print1.md}
явно сформулированы внутреннее время, микроскопический горизонт и число
$1{,}33$. Однако в архиве не найдено ни уравнений конкретной краевой задачи,
ни нормировки волновой функции, ни исходного кода, ни таблицы сходимости,
из которых это число можно воспроизвести.

Более того, безразмерная стационарная система Шрёдингера--Ньютона обладает
масштабированием. Если $(\psi(r),\Phi(r),\epsilon)$ является решением, то
при подходящем выборе единиц
\begin{equation}
 \psi_\lambda(r)=\lambda^2\psi(\lambda r),\qquad
 \Phi_\lambda(r)=\lambda^2\Phi(\lambda r),\qquad
 \epsilon_\lambda=\lambda^2\epsilon.
 \label{eq:v5-sn-scaling}
\end{equation}
При этом норма масштабируется как $M_\lambda=\lambda M$, а характерный
радиус как $R_\lambda=R/\lambda$. Поэтому отдельно взятое число радиуса не
является универсальным: требуется заранее фиксировать норму, массу и
систему единиц. Численные стационарные и эволюционные решения этой системы
обсуждаются, например, в \path{arXiv:gr-qc/0404014}; они не выделяют
архивное число проекта без дополнительной нормировки.

\begin{nogocriterion}[Невоспроизводимость исторического радиуса]
Число $R=1{,}33$ не переносится в действующую ветвь как результат. Без
сохранившейся краевой задачи и нормировки оно является историческим
сообщением, а масштабная симметрия системы Шрёдингера--Ньютона запрещает
считать его безусловным безразмерным предсказанием.
\end{nogocriterion}

\section{Минимальная переменная внутреннего времени}

Пусть $q$ измеряет амплитуду геометрического дефекта, а
$N(q)\in(0,1]$ --- коэффициент хода собственных часов:
\begin{equation}
 d\tau=N(q)\,dt.
 \label{eq:v5-lapse-definition}
\end{equation}
Здесь $t$ есть время внешнего наблюдателя, а $\tau$ --- внутреннее время.
Фраза ``замедление в четыре раза'' соответствует
\begin{equation}
 N(q_*)=\frac14.
 \label{eq:v5-quarter-lapse}
\end{equation}

Если внутренний процесс имеет скорость распада $\Gamma_{\rm int}$ на
единицу собственного времени, то внешний наблюдатель получает
\begin{equation}
 \Gamma_{\rm ext}=N(q_*)\Gamma_{\rm int},\qquad
 T_{\rm ext}=\frac{1}{N(q_*)\Gamma_{\rm int}}.
 \label{eq:v5-external-lifetime}
\end{equation}
Следовательно, $N=1/4$ увеличивает время жизни ровно в четыре раза, но не
делает объект вечным. Бесконечное внешнее время жизни потребовало бы
$N\to0$ либо собственного запрета распада $\Gamma_{\rm int}=0$.

\section{Почему замедление ещё не создаёт устойчивость}

Для проверки самосогласования рассмотрим минимальный эффективный
функционал
\begin{equation}
 \mathcal F(q;E)=V(q)+E N(q),
 \label{eq:v5-time-well-functional}
\end{equation}
где $V$ есть энергия образования дефекта, а $E$ --- внутренняя энергия
циркулирующего режима, измеренная в его собственных единицах. Ненулевое
состояние $q_*$ устойчиво только при двух независимых условиях
\begin{equation}
 V'(q_*)+E N'(q_*)=0,
 \qquad
 V''(q_*)+E N''(q_*)>0.
 \label{eq:v5-time-well-stability}
\end{equation}

Условие $N(q_*)=1/4$ не определяет ни одно из них. Это видно на явном
семействе. Положим $x=q-q_*$ и
\begin{equation}
 N(q)=\frac14 e^{-a x},\qquad
 V(q)=\frac{Ea}{4}x+\frac{\kappa}{2}x^2+
       \frac{\lambda}{4}x^4,qquad \lambda>0.
 \label{eq:v5-same-lapse-family}
\end{equation}
Для всех $\kappa$ выполняются $N(q_*)=1/4$ и $\mathcal F'(q_*)=0$, но
\begin{equation}
 \mathcal F''(q_*)=\kappa+\frac{Ea^2}{4}
 \label{eq:v5-same-lapse-hessian}
\end{equation}
может быть положительной, нулевой или отрицательной. Одинаковое
четырёхкратное замедление совместимо и с устойчивостью, и с критической
точкой, и с распадом.

\begin{theorem}[Разделение часов и устойчивости]
Коэффициент внутреннего времени фиксирует отношение скоростей процессов,
но не знак гессиана дефекта. Для получения частицы требуется независимо
вывести энергию $V$, обратную реакцию $N(q)$ и их относительную
нормировку. В текущем родителе $H_{15}/M_{35}$ такого поля хода времени и
такой связи нет.
\end{theorem}

\section{Горизонт, сингулярность и открытый обмен}

Буквальный горизонт событий является слишком сильной реализацией раннего
образа: по определению он отделяет область, из которой классический сигнал
не выходит наружу. Это противоречит требованию, чтобы свет свободно входил
в частицу и выходил из неё. Кроме того, горизонт не является
сингулярностью; смешение этих понятий в раннем тексте должно быть снято.

Более содержательной является не буквальная чёрная дыра, а
\emph{эффективный горизонт переноса}: область малого $N(q)$, которая долго
задерживает внутреннюю фазу, но не обнуляет связь полностью. Тогда объект
может быть долгоживущим резонансом. Однако в пределе $N\to0$ внешний
энергетический вклад $E_\infty=N E$ также стремится к нулю при фиксированном
$E$, а свободный обмен требует отдельного ненулевого канала связи.

Внешняя литература действительно содержит близкую идею геона ---
самогравитирующей конфигурации поля. Но гладкие периодические геоны обычно
требуют дополнительных удерживающих условий, например
антидеситтеровской границы; асимптотически плоский чисто
электромагнитный маршрут не даёт готовой устойчивой элементарной частицы.
См. историко-технический аудит \path{arXiv:1708.05674} и
антидеситтеровские конструкции \path{arXiv:1408.5906}. Поэтому
заимствование слова ``геон'' не закрывает проектный механизм.

\section{Сильное прочтение исходной интуиции}

Раннюю мысль можно сохранить без утверждения, что одна и та же порция
света навечно заперта внутри. Более устойчивый образ таков:
\begin{quote}
частица есть самоподдерживающийся геометрический дефект, через который
световой перенос может проходить; сохраняется форма потока, а не
обязательно каждый находившийся в ней квант возбуждения.
\end{quote}
Так вихрь в жидкости сохраняет узнаваемую форму, хотя составляющие его
молекулы меняются. Эта интерпретация снимает требование вечного хранения
того же сигнала, но переводит задачу в существенно нелинейную архитектуру:
поток должен создавать дефект, дефект должен направлять поток, а связанная
пара должна иметь устойчивый ненулевой минимум.

\section{Статус и критерий остановки}

Получен не тупик всей идеи, а точная граница её нынешней версии:
\begin{enumerate}
 \item внутреннее время можно формализовать коэффициентом $N(q)$;
 \item значение $N=1/4$ даёт четырёхкратную задержку, не вечность;
 \item историческое $R=1{,}33$ невоспроизводимо и не универсально;
 \item устойчивость требует отдельного положительного гессиана;
 \item буквальный горизонт несовместим со свободным выходом;
 \item эффективная временная яма остаётся допустимой как нелинейный
 долгоживущий дефект;
 \item действующий родитель не содержит необходимых полей обратной реакции.
\end{enumerate}

Продолжение внутри Тома V останавливается. Добавление динамической метрики,
коэффициента $N(q)$ и самогравитирующего поля было бы новой родительской
архитектурой, а не ещё одним уточнением текущего $H_{15}/M_{35}$. Следующей
главе надлежит подвести итог части II и сформулировать условия возможной
версии VI: нелинейный совместный функционал ``поток + дефект + ход
времени'', заранее нормированный и проверяемый на устойчивость и открытый
обмен.

\begin{protocol}
\begin{enumerate}
 \item ранняя гипотеза внутреннего времени найдена: \textbf{да};
 \item исходный расчёт $R=1{,}33$ найден: \textbf{нет};
 \item универсальность $R=1{,}33$: \textbf{опровергнута без нормировки};
 \item коэффициент $N=1/4$: \textbf{формализован};
 \item вечность из $N=1/4$: \textbf{нет, только множитель четыре};
 \item устойчивость из одного замедления: \textbf{не следует};
 \item буквальный горизонт со свободным выходом: \textbf{несовместим};
 \item нелинейный дефект с проходящим потоком: \textbf{допустимая новая
 архитектура};
 \item физическое замыкание: \textbf{не достигнуто}.
\end{enumerate}
\end{protocol}

Машинный протокол сохранён в
\path{s2t_v5_self_consistent_internal_time_horizon_gate_results.json}.